\documentclass{article}
\usepackage{fleqn}
\usepackage{epsf}
\usepackage{aima2e-slides}

\begin{document}

\begin{huge}
\titleslide{Suy diễn trong mạng Bayes (Inference in Bayesian networks)}{Chương 14.4--5}

\sf

%%%%%%%%%%%% Slide %%%%%%%%%%%%%%%%%%%%%%%%%%%%%%%%%%%%%%%%%%%%%%%%%%%
\heading{Phác thảo}

\blob Suy luận chính xác bằng cách liệt kê

\blob Suy luận chính xác bằng cách loại bỏ biến

\blob Suy luận gần đúng bằng mô phỏng ngẫu nhiên

\blob Suy luận gần đúng của chuỗi Markov Monte Carlo


%%%%%%%%%%%% Slide %%%%%%%%%%%%%%%%%%%%%%%%%%%%%%%%%%%%%%%%%%%%%%%%%%%
\heading{Nhiệm vụ suy luận}

\defn{Truy vấn đơn giản}: tính biên sau \mat{$\pv(X_i|\mbf{E}\eq \e)$}\al
  ví dụ: \mat{$P(NoGas|Gauge\eq empty,Lights\eq on,Starts\eq false)$}

\defn{Truy vấn liên kết}: \mat{$\pv(X_i,X_j|\mbf{E}\eq \e) = 
    \pv(X_i|\mbf{E}\eq \e)  \pv(X_j|X_i,\mbf{E}\eq \e) $}

\defn{Quyết định tối ưu}: mạng quyết định bao gồm thông tin tiện ích;\nl
    suy luận xác suất cần thiết cho \mat{$P(outcome|action,evidence)$}

\defn{Giá trị của thông tin}: cần tìm kiếm bằng chứng nào tiếp theo?

\defn{Phân tích độ nhạy}: giá trị xác suất nào là quan trọng nhất?

\defn{Giải thích}: tại sao tôi cần động cơ khởi động mới?



%%%%%%%%%%%% Slide %%%%%%%%%%%%%%%%%%%%%%%%%%%%%%%%%%%%%%%%%%%%%%%%%%%
\heading{Suy luận bằng phép liệt kê}

Cách hơi thông minh để tính tổng các biến từ khớp
mà không thực sự xây dựng sự biểu diễn rõ ràng của nó

Truy vấn đơn giản trên mạng trộm:\hspace*{2.5in}\epsfxsize=1.5in\raisebox{-1.5in[0pt][0pt]{\epsffile{\file{figures}{burglary-small.ps}}\\
\mat{$\pv(B|j,m)$}\\
\mat{$= \pv(B,j,m)/P(j,m)$}\\
\mat{$= \alpha \pv(B,j,m)$}\\
\mat{$= \alpha \ \mysum_e \ \mysum_a \ \pv(B,e,a,j,m)$}

Viết lại toàn bộ các mục chung sử dụng tích của các mục CPT:\\
\mat{$\pv(B|j,m)$}\\
\mat{$= \alpha \ \mysum_e \ \mysum_a \ \pv(B)P(e)\pv(a|B,e)P(j|a)P(m|a) $}\\
\mat{$= \alpha \pv(B)\ \mysum_e\ P(e)\ \mysum_a\ \pv(a|B,e)P(j|a)P(m|a)$}
       
Đệ quy liệt kê theo chiều sâu đầu tiên: \mat{$O(n)$} không gian, \mat{$O(d^n)$} thời gian


%%%%%%%%%%%% Slide %%%%%%%%%%%%%%%%%%%%%%%%%%%%%%%%%%%%%%%%%%%%%%%%%%%
\heading{Thuật toán liệt kê}

\input{\file{algorithms}{enumeration-algorithm}}

%%%%%%%%%%%% Slide %%%%%%%%%%%%%%%%%%%%%%%%%%%%%%%%%%%%%%%%%%%%%%%%%%%
\heading{Cây đánh giá}

\epsfxsize=1.0\textwidth
\fig{\file{figures}{enumeration-tree.ps}}

Việc liệt kê không hiệu quả: tính toán lặp đi lặp lại \al
  ví dụ: tính \mat{$P(j|a)P(m|a)$} cho mỗi giá trị của \mat{$e$}




%%%%%%%%%%%% Slide %%%%%%%%%%%%%%%%%%%%%%%%%%%%%%%%%%%%%%%%%%%%%%%%%%%
\heading{Suy luận bằng cách loại bỏ biến}

Loại bỏ biến: thực hiện tính tổng từ phải sang trái,\\
lưu trữ kết quả trung gian (\defn{factors}) để tránh tính toán lại

\mat{$\pv(B|j,m)$}\nl
    \mat{$= \alpha \underbrace{\pv(B)}_B 
             \mysum_e \underbrace{P(e)}_E
             \mysum_a \underbrace{\pv(a|B,e)}_A
             \underbrace{P(j|a)}_J
             \underbrace{P(m|a)}_M$}\nl
    \mat{$= \alpha \pv(B)\mysum_e P(e)\mysum_a \pv(a|B,e) P(j|a) f_M(a)$}\nl
    \mat{$= \alpha \pv(B)\mysum_e P(e)\mysum_a \pv(a|B,e) f_J(a) f_M(a)$}\nl
    \mat{$= \alpha \pv(B)\mysum_e P(e)\mysum_a f_A(a,b,e) f_J(a) f_M(a)$}\nl
    \mat{$= \alpha \pv(B)\mysum_e P(e)f_{\bar{A}JM}(b,e) $} (tổng \mat{$A$})\nl
    \mat{$= \alpha \pv(B)f_{\bar{E}\bar{A}JM}(b)$} (tổng \mat{$E$})\nl
    \mat{$= \alpha f_B(b)\stimes f_{\bar{E}\bar{A}JM}(b)$} 


%%%%%%%%%%%% Slide %%%%%%%%%%%%%%%%%%%%%%%%%%%%%%%%%%%%%%%%%%%%%%%%%%%
\heading{Loại bỏ biến: Các thao tác cơ bản}

\defn{Tính tổng} một biến từ tích của các thừa số: \al
  di chuyển bất kỳ yếu tố không đổi nào ra ngoài phép tính tổng \al
  cộng các ma trận con theo tích từng điểm của các thừa số còn lại

\mat{$\mysum_x f_1 \stimes \cdots \stimes f_k =
f_1 \stimes \cdots \stimes f_i\  \mysum_x\; f_{i+1} \stimes \cdots \stimes
f_k = f_1 \stimes \cdots \stimes f_i \stimes f_{\bar{X}}$}

giả sử \mat{$f_1,\ldots,f_i$} không phụ thuộc vào \mat{$X$}


\defn{Tích điểm} của thừa số \mat{$f_1$} và \mat{$f_2$}:\al
  \mat{$f_1(x_1,\ldots,x_j,y_1,\ldots,y_k) \stimes
     f_2(y_1,\ldots,y_k,z_1,\ldots,z_l)$}\nl
  = \mat{$f(x_1,\ldots,x_j,y_1,\ldots,y_k,z_1,\ldots,z_l)$}\\
Ví dụ: \mat{$f_1(a,b) \stimes f_2(b,c) = f(a,b,c)$}





%%%%%%%%%%%% Slide %%%%%%%%%%%%%%%%%%%%%%%%%%%%%%%%%%%%%%%%%%%%%%%%%%%
\heading{Thuật toán loại bỏ biến }

\input{\file{algorithms}{elimination-ask-algorithm}}


%%%%%%%%%%%% Slide %%%%%%%%%%%%%%%%%%%%%%%%%%%%%%%%%%%%%%%%%%%%%%%%%%%
\heading{Các biến không liên quan}

Hãy xem xét truy vấn \mat{$P(JohnCalls|Burglary\eq true)$}\hspace*{1.0in}\epsfxsize=1.5in\raisebox{-1.5in[0pt][0pt]{\epsffile{\file{figures}{burglary-small.ps}}}
\mat{\[
  P(J|b) = \alpha P(b) \sum_e P(e) \sum_a P(a|b,e) P(J|a) \sum_m P(m|a)
\]}
Tổng trên \mat{$m$} giống hệt 1; \mat{$M$} là \emph{không liên quan} với truy vấn

\vspace*{0.3in}

Thm 1: \mat{$Y$} không liên quan trừ khi \mat{$Y\elt Ancestors(\{X\}\union \E)$}

Đây, \mat{$X\eq JohnCalls$}, \mat{$\E\eq\{Burglary\}$} và \\
\mat{$Ancestors(\{X\}\union \E) = \{Alarm,Earthquake\}$}\\
vì vậy \mat{$MaryCalls$} không liên quan

(So sánh điều này với chuỗi ngược từ truy vấn trong KB mệnh đề Horn)

%%%%%%%%%%%% Slide %%%%%%%%%%%%%%%%%%%%%%%%%%%%%%%%%%%%%%%%%%%%%%%%%%%
\heading{Các biến không liên quan tiếp.}

Defn: \underline{biểu đồ đạo đức} của Bayes net: kết hôn với tất cả cha mẹ và mũi tên thả

Định nghĩa: \mat{$\A$} được phân tách bằng \underline{m} khỏi \mat{$\B$} bởi \mat{$\C$} nếu được phân tách bằng \mat{$\C$} trong biểu đồ đạo đức

Thm 2: \mat{$Y$} không liên quan nếu m được phân tách khỏi \mat{$X$} bởi \mat{$\E$}\hspace*{1.0in}\epsfxsize=1.5in\raisebox{-1.5in[0pt][0pt]{\epsffile{\file{figures}{burglary-moral.ps}}}

\vspace*{0.3in}

Đối với \mat{$P(JohnCalls|Alarm\eq true)$}, cả \\
\mat{$Burglary$} và \mat{$Earthquake$} không liên quan

%%%%%%%%%%%% Slide %%%%%%%%%%%%%%%%%%%%%%%%%%%%%%%%%%%%%%%%%%%%%%%%%%%
\heading{Độ phức tạp của suy luận chính xác}

\defn{Mạng được kết nối đơn} (hoặc \defn{polytrees}):\al
  -- bất kỳ hai nút nào cũng được kết nối bằng nhiều nhất một đường dẫn (vô hướng)\al
  -- chi phí về thời gian và không gian của việc loại bỏ biến đổi là \mat{$O(d^k n)$}

\defn{Các mạng được kết nối nhiều lần}:\al
  -- có thể giảm 3SAT để suy luận chính xác \mat{$\implies$} NP-hard\al
  -- tương đương với \emph{đếm} kiểu 3SAT \mat{$\implies$} \#P-complete

\vspace*{0.2in}

\epsfxsize=0,75\textwidth
\fig{\file{figures}{bn-3sat.ps}}


%%%%%%%%%%%% Slide %%%%%%%%%%%%%%%%%%%%%%%%%%%%%%%%%%%%%%%%%%%%%%%%%%%
\heading{Suy luận bằng mô phỏng ngẫu nhiên}

Ý tưởng cơ bản:\al
  1) Vẽ \mat{$N$} mẫu từ phân phối lấy mẫu \mat{$S$}\hspace*{1.5in}\epsfxsize=1.0in\raisebox{-1.5in[0pt][0pt]{\epsffile{\file{figures}{coin-flip.ps}}\al
  2) Tính xác suất hậu nghiệm gần đúng \mat{$\hat P$}\al
  3) Chứng tỏ điều này hội tụ về xác suất đúng \mat{$P$}

Tóm tắt:\al
  -- Lấy mẫu từ mạng trống\al
  -- Lấy mẫu từ chối: từ chối các mẫu không đồng ý với bằng chứng\al
  -- Trọng số khả năng: sử dụng bằng chứng để cân mẫu\al
  -- Chuỗi Markov Monte Carlo (MCMC): mẫu từ quá trình ngẫu nhiên\nl
      có phân bố cố định là phần sau thực sự


%%%%%%%%%%%% Slide %%%%%%%%%%%%%%%%%%%%%%%%%%%%%%%%%%%%%%%%%%%%%%%%%%%
\heading{Lấy mẫu từ mạng trống}

\input{\file{algorithms}{prior-sample-algorithm}}

%%%%%%%%%%%% Slide %%%%%%%%%%%%%%%%%%%%%%%%%%%%%%%%%%%%%%%%%%%%%%%%%%%
\heading{Ví dụ}

\epsfxsize=0,85\textwidth
\fig{\sfile{figures}{rain-prior-sample1.ps}}

%%%%%%%%%%%% Slide %%%%%%%%%%%%%%%%%%%%%%%%%%%%%%%%%%%%%%%%%%%%%%%%%%%
\heading{Ví dụ}

\epsfxsize=0,85\textwidth
\fig{\sfile{figures}{rain-prior-sample2.ps}}

%%%%%%%%%%%% Slide %%%%%%%%%%%%%%%%%%%%%%%%%%%%%%%%%%%%%%%%%%%%%%%%%%%
\heading{Ví dụ}

\epsfxsize=0,85\textwidth
\fig{\sfile{figures}{rain-prior-sample3.ps}}

%%%%%%%%%%%% Slide %%%%%%%%%%%%%%%%%%%%%%%%%%%%%%%%%%%%%%%%%%%%%%%%%%%
\heading{Ví dụ}

\epsfxsize=0,85\textwidth
\fig{\sfile{figures}{rain-prior-sample4.ps}}

%%%%%%%%%%%% Slide %%%%%%%%%%%%%%%%%%%%%%%%%%%%%%%%%%%%%%%%%%%%%%%%%%%
\heading{Ví dụ}

\epsfxsize=0,85\textwidth
\fig{\sfile{figures}{rain-prior-sample5.ps}}

%%%%%%%%%%%% Slide %%%%%%%%%%%%%%%%%%%%%%%%%%%%%%%%%%%%%%%%%%%%%%%%%%%
\heading{Ví dụ}

\epsfxsize=0,85\textwidth
\fig{\sfile{figures}{rain-prior-sample6.ps}}

%%%%%%%%%%%% Slide %%%%%%%%%%%%%%%%%%%%%%%%%%%%%%%%%%%%%%%%%%%%%%%%%%%
\heading{Ví dụ}

\epsfxsize=0,85\textwidth
\fig{\sfile{figures}{rain-prior-sample7.ps}}



%%%%%%%%%%%% Slide %%%%%%%%%%%%%%%%%%%%%%%%%%%%%%%%%%%%%%%%%%%%%%%%%%%
\heading{Lấy mẫu từ mạng trống tiếp theo.}

Xác suất \prog{PriorSample} tạo ra một sự kiện cụ thể\al
  \mat{$S_{PS}(x_1\ldots x_n) = \myprod_{i\eq 1}^n P(x_i | \parents(X_i))
    = P(x_1\ldots x_n)$}\\
tức là xác suất trước thực sự

Ví dụ: \mat{$S_{PS}(t, f, t, t) = 0.5 \stimes 0.9 \stimes 0.8 \stimes 0.9 = 0.324 = P(t, f, t,t)$}

Gọi \mat{$N_{PS}(x_1\ldots x_n)$} là số lượng mẫu được tạo
cho sự kiện \mat{$x_1,\ldots, x_n$}

Sau đó chúng tôi có
\mat{\begin{eqnarray*}
  \lim_{N\to\infty} \hat P(x_1,\ldots, x_n) 
      & = & \lim_{N\to\infty} N_{PS}(x_1,\ldots, x_n)/N \\
      & = & S_{PS}(x_1,\ldots,x_n) \\
      & = & P(x_1\ldots x_n)
\end{eqnarray*}}
Nghĩa là, các ước tính bắt nguồn từ \prog{PriorSample} là \defn{nhất quán}

Viết tắt: \mat{$\hat P(x_1,\ldots, x_n) \approx P(x_1\ldots x_n)$}



%%%%%%%%%%%% Slide %%%%%%%%%%%%%%%%%%%%%%%%%%%%%%%%%%%%%%%%%%%%%%%%%%%
\heading{Lấy mẫu loại bỏ}

\mat{$\hat{\pv}(X|\e)$} được ước tính từ các mẫu phù hợp với \mat{$\e$}

\input{\file{algorithms}{rejection-sampling-algorithm}}

Ví dụ: ước tính \mat{$\pv(Rain|Sprinkler\eq true)$} sử dụng 100 mẫu \al
  27 mẫu có \mat{$Sprinkler\eq true$}\nl
    Trong số này, 8 có \mat{$Rain\eq true$} và 19 có \mat{$Rain\eq false$}.\\[0.2in]
\mat{$\hat{\pv}(Rain|Sprinkler\eq true) = \noprog{Normalize}(\<8,19\>) =
\<0.296,0.704\>$}

Tương tự như quy trình ước lượng thực nghiệm cơ bản trong thế giới thực

%%%%%%%%%%%% Slide %%%%%%%%%%%%%%%%%%%%%%%%%%%%%%%%%%%%%%%%%%%%%%%%%%%
\heading{Phân tích lấy mẫu từ chối}


\mat{$\hat{\pv}(X|\e) = \alpha \mbf{N}_{PS}(X,\e) $} 
   \qquad (định nghĩa thuật toán)\al
\mat{$= \mbf{N}_{PS}(X,\e)/N_{PS}(\e)$} 
   \qquad (chuẩn hóa bởi \mat{$N_{PS}(\e)$})\al
\mat{$\approx \pv(X,\e)/P(\e)$} 
   \qquad (thuộc tính của \prog{PriorSample})\al
\mat{$= \pv(X|\e)$} 
   \qquad (định nghĩa về xác suất có điều kiện)

Do đó việc lấy mẫu từ chối trả về các ước tính sau nhất quán

Vấn đề: cực kỳ tốn kém nếu \mat{$P(\e)$} nhỏ

\mat{$P(\e)$} giảm theo cấp số nhân với số lượng biến bằng chứng!


%%%%%%%%%%%% Slide %%%%%%%%%%%%%%%%%%%%%%%%%%%%%%%%%%%%%%%%%%%%%%%%%%%
\heading{Trọng số khả năng}

Ý tưởng: sửa các biến bằng chứng, chỉ lấy mẫu các biến không có bằng chứng,\\
và cân nhắc từng mẫu theo khả năng nó phù hợp với bằng chứng

\input{\file{algorithms}{likelihood-weighting-algorithm}}


%%%%%%%%%%%% Slide %%%%%%%%%%%%%%%%%%%%%%%%%%%%%%%%%%%%%%%%%%%%%%%%%%%
\heading{Ví dụ về trọng số khả năng}

\epsfxsize=0,7\textwidth
\fig{\sfile{figures}{rain-lw-sample1.ps}}

\mat{$w = 1.0$}

%%%%%%%%%%%% Slide %%%%%%%%%%%%%%%%%%%%%%%%%%%%%%%%%%%%%%%%%%%%%%%%%%%
\heading{Ví dụ về trọng số khả năng}

\epsfxsize=0,7\textwidth
\fig{\sfile{figures}{rain-lw-sample2.ps}}

\mat{$w = 1.0$}

%%%%%%%%%%%% Slide %%%%%%%%%%%%%%%%%%%%%%%%%%%%%%%%%%%%%%%%%%%%%%%%%%%
\heading{Ví dụ về trọng số khả năng}

\epsfxsize=0,7\textwidth
\fig{\sfile{figures}{rain-lw-sample3.ps}}

\mat{$w = 1.0$}

%%%%%%%%%%%% Slide %%%%%%%%%%%%%%%%%%%%%%%%%%%%%%%%%%%%%%%%%%%%%%%%%%%
\heading{Ví dụ về trọng số khả năng}

\epsfxsize=0,7\textwidth
\fig{\sfile{figures}{rain-lw-sample3.ps}}

\mat{$w = 1.0 \stimes 0.1$}

%%%%%%%%%%%% Slide %%%%%%%%%%%%%%%%%%%%%%%%%%%%%%%%%%%%%%%%%%%%%%%%%%%
\heading{Ví dụ về trọng số khả năng}

\epsfxsize=0,7\textwidth
\fig{\sfile{figures}{rain-lw-sample4.ps}}

\mat{$w = 1.0 \stimes 0.1$}

%%%%%%%%%%%% Slide %%%%%%%%%%%%%%%%%%%%%%%%%%%%%%%%%%%%%%%%%%%%%%%%%%%
\heading{Ví dụ về trọng số khả năng}

\epsfxsize=0,7\textwidth
\fig{\sfile{figures}{rain-lw-sample5.ps}}

\mat{$w = 1.0 \stimes 0.1$}

%%%%%%%%%%%% Slide %%%%%%%%%%%%%%%%%%%%%%%%%%%%%%%%%%%%%%%%%%%%%%%%%%%
\heading{Ví dụ về trọng số khả năng}

\epsfxsize=0,7\textwidth
\fig{\sfile{figures}{rain-lw-sample5.ps}}

\mat{$w = 1.0 \stimes 0.1 \stimes 0.99 = 0.099$}

%%%%%%%%%%%% Slide %%%%%%%%%%%%%%%%%%%%%%%%%%%%%%%%%%%%%%%%%%%%%%%%%%%
\heading{Phân tích trọng số khả năng}

Xác suất lấy mẫu cho \prog{WeightedSample} là\al
  \mat{$S_{WS}(\mbf{z},\e) = \myprod_{i\eq 1}^l P(z_i|\parents(Z_i))$}\\
Lưu ý: chỉ chú ý đến bằng chứng trong \emph{tổ tiên}\hspace*{0.3in}\epsfxsize=2.1in\raisebox{-1.3in[0pt][0pt]{\epsffile{\file{figures}{rain-lw2.ps}}\nl
    \mat{$\implies$} đâu đó ``ở giữa'' trước và\nl
    phân bố sau

Trọng lượng của một mẫu nhất định \mat{$\mbf{z},\e$} là \al
  \mat{$w(\mbf{z},\e) = \myprod_{i\eq 1}^m P(e_i | \parents(E_i))$}

Xác suất lấy mẫu có trọng số là \al
  \mat{$S_{WS}(\mbf{z},\e) w(\mbf{z},\e)$}\nl
    \mat{$= \myprod_{i\eq 1}^l P(z_i|\parents(Z_i))\ \ 
       \myprod_{i\eq 1}^m P(e_i|\parents(E_i))$}\nl
    \mat{$= P(\mbf{z},\e)$} (theo ngữ nghĩa toàn cầu tiêu chuẩn của mạng)

Do đó, trọng số khả năng trả về ước tính nhất quán\\
nhưng hiệu suất vẫn suy giảm với nhiều biến số bằng chứng\\
vì một số mẫu có gần như toàn bộ trọng lượng

%%%%%%%%%%%% Slide %%%%%%%%%%%%%%%%%%%%%%%%%%%%%%%%%%%%%%%%%%%%%%%%%%%
\heading{Suy luận gần đúng sử dụng MCMC}

``Trạng thái'' của mạng = sự gán hiện tại cho tất cả các biến.

Tạo trạng thái tiếp theo bằng cách lấy mẫu một biến cho chăn Markov \\
Lấy mẫu lần lượt từng biến, giữ bằng chứng cố định

\input{\file{algorithms}{mcmc-algorithm}}

Cũng có thể chọn một biến để lấy mẫu ngẫu nhiên mỗi lần

%%%%%%%%%%%% Slide %%%%%%%%%%%%%%%%%%%%%%%%%%%%%%%%%%%%%%%%%%%%%%%%%%%
\heading{Chuỗi Markov}

Với \mat{$Sprinkler\eq true,WetGrass\eq true$}, có bốn trạng thái:

\vspace*{0.2in}

\epsfxsize=0,7\textwidth
\fig{\file{figures}{rain-chain.ps}}

Đi loanh quanh một lúc, tính trung bình những gì bạn nhìn thấy


%%%%%%%%%%%% Slide %%%%%%%%%%%%%%%%%%%%%%%%%%%%%%%%%%%%%%%%%%%%%%%%%%%
\heading{Ví dụ MCMC tiếp theo.}

Ước tính \mat{$\pv(Rain|Sprinkler\eq true,WetGrass\eq true)$}

Mẫu \mat{$Cloudy$} hoặc \mat{$Rain$} được cung cấp chăn Markov, lặp lại.\\
Đếm số lần \mat{$Rain$} đúng và sai trong các mẫu.

Ví dụ: truy cập 100 tiểu bang\al
  31 có \mat{$Rain\eq true$}, 69 có \mat{$Rain\eq false$}

\mat{$\hat{\pv}(Rain|Sprinkler\eq true,WetGrass\eq true)$}\al
      \mat{$= \noprog{Normalize}(\<31,69\>) = \<0.31,0.69\>$}

Định lý: phương pháp tiếp cận chuỗi \defn{phân phối cố định}: \al
Phần thời gian dài hạn dành cho mỗi trạng thái chính xác là \al
tỷ lệ thuận với xác suất sau của nó

%%%%%%%%%%%% Slide %%%%%%%%%%%%%%%%%%%%%%%%%%%%%%%%%%%%%%%%%%%%%%%%%%%
\heading{Lấy mẫu chăn Markov}

Chăn Markov của \mat{$Cloudy$} là\hspace*{2.3in}\epsfxsize=2.1in\raisebox{-1.3in[0pt][0pt]{\epsffile{\file{figures}{rain-lw1.ps}}\nl
    \mat{$Sprinkler$} và \mat{$Rain$}\\
Chăn Markov của \mat{$Rain$} là \nl
    \mat{$Cloudy$}, \mat{$Sprinkler$} và \mat{$WetGrass$}

Xác suất cho chăn Markov được tính như sau:\al
  \mat{$P(x_i'|\markovBlanket(X_i)) = P(x_i'|\parents(X_i)) 
           \myprod_{Z_j\in Children(X_i)} P(z_j|\parents(Z_j))$}

Dễ dàng triển khai trong các hệ thống song song truyền thông điệp, bộ não

Các vấn đề tính toán chính:\al
  1) Khó biết liệu đã đạt được sự hội tụ hay chưa\al
  2) Có thể lãng phí nếu chăn Markov lớn:\nl
    \mat{$P(X_i|\markovBlanket(X_i))$} sẽ không thay đổi nhiều (luật số lớn)

%%%%%%%%%%%% Slide %%%%%%%%%%%%%%%%%%%%%%%%%%%%%%%%%%%%%%%%%%%%%%%%%%%
\heading{Tóm tắt}

Suy luận chính xác bằng cách loại bỏ biến: \al
 -- polytime trên polytrees, NP-hard trên đồ thị chung\al
 -- không gian = thời gian, rất nhạy cảm với cấu trúc liên kết

Suy luận gần đúng của LW, MCMC:\al
 -- LW hoạt động kém khi có nhiều bằng chứng (hạ lưu)\al
 -- LW, MCMC thường không nhạy cảm với cấu trúc liên kết\al
 -- Sự hội tụ có thể rất chậm với xác suất gần bằng 1 hoặc 0\al
 -- Có thể xử lý sự kết hợp tùy ý của các biến rời rạc và liên tục



\end{huge} 
\end{document}

%%%%%%%%%%%% Slide %%%%%%%%%%%%%%%%%%%%%%%%%%%%%%%%%%%%%%%%%%%%%%%%%%%
\heading{Phân tích MCMC: Đề cương}

Xác suất chuyển tiếp \mat{$\transition{\x}{\x'}$}

Xác suất chiếm dụng \mat{$\pi_t(\x)$} tại thời điểm \mat{$t$}

Điều kiện cân bằng trên \mat{$\pi_t$} xác định phân bố cố định \mat{$\pi(\x)$}\nl
Lưu ý: phân phối cố định phụ thuộc vào sự lựa chọn \mat{$\transition{\x}{\x'}$}

Theo cặp \defn{cân bằng chi tiết} ở các trạng thái đảm bảo trạng thái cân bằng

\defn{Xác suất chuyển tiếp lấy mẫu Gibbs}:\nl
    lấy mẫu từng biến cho các giá trị hiện tại của tất cả các biến khác\\
\mat{$\implies$} cân bằng chi tiết với mặt sau thật

Đối với mạng Bayesian, việc lấy mẫu Gibbs giảm xuống còn \\
lấy mẫu có điều kiện trên chăn Markov của mỗi biến


%%%%%%%%%%%% Slide %%%%%%%%%%%%%%%%%%%%%%%%%%%%%%%%%%%%%%%%%%%%%%%%%%%
\heading{Phân phối cố định}

\mat{$\pi_t(\x)$} = xác suất ở trạng thái \mat{$\x$} tại thời điểm \mat{$t$}\\
\mat{$\pi_{t+1}(\x')$} = xác suất ở trạng thái \mat{$\x'$} tại thời điểm \mat{$t+1$}

\mat{$\pi_{t+1}$} xét về \mat{$\pi_t$} và \mat{$\transition{\x}{\x'}$}
\mat{\[
  \pi_{t+1}(\x') = \mysum_{\smbf{x}} \pi_t(\x) \transition{\x}{\x'}
\]}
Phân bố cố định: \mat{$\pi_t = \pi_{t+1} = \pi$}
\mat{\[
  \pi(\x') = \mysum_{\smbf{x}} \pi(\x) \transition{\x}{\x'}
      \qquad\mbox{for all }\x'
\]}
Nếu \mat{$\pi$} tồn tại thì nó là duy nhất (cụ thể cho \mat{$\transition{\x}{\x'}$})

Ở trạng thái cân bằng, ``dòng ra'' dự kiến = ``dòng vào'' dự kiến
  

%%%%%%%%%%%% Slide %%%%%%%%%%%%%%%%%%%%%%%%%%%%%%%%%%%%%%%%%%%%%%%%%%%
\heading{Số dư chi tiết}

``Outflow'' = ``inflow'' cho mỗi cặp trạng thái:
\mat{\[
  \pi(\x) \transition{\x}{\x'} 
   = \pi(\x') \transition{\x'}{\x}
      \qquad\mbox{for all }\x,\ \x'
\]}
Cân bằng chi tiết \mat{$\implies$} tính ổn định:
\mat{\begin{eqnarray*}
\mysum_{\smbf{x}} \pi(\x) \transition{\x}{\x'} 
   & = & \mysum_{\smbf{x}} \pi(\x') \transition{\x'}{\x} \\
   & = & \pi(\x') \mysum_{\smbf{x}} \transition{\x'}{\x} \\
   & = & \pi(\x')
\end{eqnarray*}}

Các thuật toán MCMC thường được xây dựng bằng cách thiết kế một quá trình chuyển đổi\\
xác suất \mat{$q$} cân bằng chi tiết với \mat{$\pi$} mong muốn


%%%%%%%%%%%% Slide %%%%%%%%%%%%%%%%%%%%%%%%%%%%%%%%%%%%%%%%%%%%%%%%%%%
\heading{ Lấy mẫu Gibbs }

Lấy mẫu lần lượt từng biến, cho trước \emph{tất cả các biến khác}

Lấy mẫu \mat{$X_i$}, đặt \mat{$\bar{\X_i}$} là tất cả các biến không có bằng chứng khác\\
Giá trị hiện tại là \mat{$x_i$} và \mat{$\bar{\x_i}$}; \mat{$\e$} đã được sửa \\
Xác suất chuyển tiếp được đưa ra bởi
\mat{\[
\transition{\x}{\x'} =
  \transition{x_i,\bar{\x_i}}{x_i',\bar{\x_i}} =
    P(x_i'|\bar{\x_i},\e)
\]}
Điều này mang lại sự cân bằng chi tiết với phần sau thực sự \mat{$P(\x|\e)$}:
\mat{\begin{eqnarray*}
\pi(\x) \transition{\x}{\x'} 
   &=& P(\x|\e) P(x_i'|\bar{\x_i},\e) 
     =  P(x_i,\bar{\x_i}|\e)P(x_i'|\bar{\x_i},\e)  \\
   &=& P(x_i|\bar{\x_i},\e)P(\bar{\x_i}|\e)
       P(x_i'|\bar{\x_i},\e) \quad\mbox{(chain rule)} \\
   &=& P(x_i|\bar{\x_i},\e)P(x_i',\bar{\x_i}|\e)
      \qquad\mbox{(chain rule backwards)} \\
   &=& \transition{\x'}{\x} \pi(\x') 
       = \pi(\x') \transition{\x'}{\x} 
\end{eqnarray*}}

%%%%%%%%%%%% Slide %%%%%%%%%%%%%%%%%%%%%%%%%%%%%%%%%%%%%%%%%%%%%%%%%%%
\heading{Hiệu suất của thuật toán xấp xỉ}

\defn{Xấp xỉ tuyệt đối}: 
  \mat{$|P(X|\e) - \hat P(X|\e)| \leq \epsilon$}

\defn{Xấp xỉ tương đối}: 
  \mat{$\frac{|P(X|\smbf{e}) - \hat P(X|\smbf{e})|}{P(X|\smbf{e})} \leq \epsilon$}

Tương đối \mat{$\implies$} tuyệt đối kể từ \mat{$0\leq P \leq 1$} (có thể là \mat{$O(2^{-n})$})

Các thuật toán ngẫu nhiên có thể thất bại với xác suất tối đa \mat{$\delta$}

Xấp xỉ đa thời gian: \mat{$\mbox{poly}(n,\epsilon^{-1},\log \delta^{-1})$}

Định lý (Dagum và Luby, 1993): cả tuyệt đối và tương đối\\
xấp xỉ cho các thuật toán xác định hoặc ngẫu nhiên\\
là NP-hard cho mọi \mat{$\epsilon,\delta<0.5$}

(Polytime gần đúng tuyệt đối không có bằng chứng---giới hạn Chernoff)

